\section*{Problem 1}

Given the reads
\[
    \texttt{\$ATGC},\ \texttt{ATGC},\ \texttt{TGCAT},\ \texttt{ATGCT},\ \texttt{GCTAG},\ \texttt{CTAGC},\ \texttt{TAGCA},\ \texttt{AGCA\$},\ \texttt{ATGCA},\ \texttt{GCATG},\ \texttt{TGCTA},\ \texttt{CATGC},
\]
construct the de~Bruijn graph with \emph{4-mers as nodes} (so edges are 5-mers) and output the assembled sequence.

\paragraph{Method (concise).}
With \(k=4\)-mer nodes, each observed 5-mer becomes a directed edge from its 4-mer prefix to its 4-mer suffix. The contig is recovered by an Eulerian path that traverses each edge once; spelling the path yields the assembled sequence.

\paragraph{Step 1: 5-mers (edges).}
Sliding a window of length 5 over all reads and deduplicating yields the set of edges:
\[
    \{\ \texttt{\$ATGC},\ \texttt{TGCAT},\ \texttt{ATGCT},\ \texttt{GCTAG},\ \texttt{CTAGC},\ \texttt{TAGCA},\ \texttt{AGCA\$},\ \texttt{ATGCA},\ \texttt{GCATG},\ \texttt{TGCTA},\ \texttt{CATGC}\ \}.
\]

\paragraph{Step 2: Edge mapping (prefix~\(\to\)~suffix).}
\[
    \begin{aligned}
        \texttt{\$ATGC} & :~\texttt{\$ATG}\to\texttt{ATGC} &
        \texttt{TGCAT}  & :~\texttt{TGCA}\to\texttt{GCAT}  &
        \texttt{ATGCT}  & :~\texttt{ATGC}\to\texttt{TGCT}  &
        \texttt{GCTAG}  & :~\texttt{GCTA}\to\texttt{CTAG}        \\
        \texttt{CTAGC}  & :~\texttt{CTAG}\to\texttt{TAGC}  &
        \texttt{TAGCA}  & :~\texttt{TAGC}\to\texttt{AGCA}  &
        \texttt{AGCA\$} & :~\texttt{AGCA}\to\texttt{GCA\$} &
        \texttt{ATGCA}  & :~\texttt{ATGC}\to\texttt{TGCA}        \\
        \texttt{GCATG}  & :~\texttt{GCAT}\to\texttt{CATG}  &
        \texttt{TGCTA}  & :~\texttt{TGCT}\to\texttt{GCTA}  &
        \texttt{CATGC}  & :~\texttt{CATG}\to\texttt{ATGC}  &   &
    \end{aligned}
\]

\paragraph{Step 3: Degree check and Eulerian path.}
All nodes are balanced except:
\[
    \texttt{\$ATG}:\ \text{out}=1,\ \text{in}=0\quad(\text{start}),\qquad
    \texttt{GCA\$}:\ \text{in}=1,\ \text{out}=0\quad(\text{end}).
\]
Hence an Eulerian path exists from \texttt{\$ATG} to \texttt{GCA\$}. One valid traversal is:
\[
    \texttt{\$ATG}\to\texttt{ATGC}\to\texttt{TGCA}\to\texttt{GCAT}\to\texttt{CATG}\to\texttt{ATGC}\to\texttt{TGCT}\to\texttt{GCTA}\to\texttt{CTAG}\to\texttt{TAGC}\to\texttt{AGCA}\to\texttt{GCA\$}.
\]

\paragraph{Step 4: Assembled sequence (spell the path).}
Start with the first 4-mer \texttt{\$ATG}; each step appends the last character of the next node. The result is
\[
    \boxed{\texttt{\$ATGCATGCTAGCA\$}}
\]

\bigskip

\begin{center}
    \resizebox{\textwidth}{!}{%  % 自动调整到页面宽度
        \begin{tikzpicture}[
                >=Latex,
                node distance=1.7cm and 1.8cm,
                every node/.style={font=\ttfamily, draw, rounded corners=2pt, inner sep=3pt},
                edge label/.style={font=\ttfamily\small, sloped, above},
                bend/.style={bend left=12},
                bendrev/.style={bend right=14}
            ]
            % Nodes (4-mers)
            \node (START) {\$ATG};
            \node (ATGC) [right=of START] {ATGC};
            \node (TGCT) [right=of ATGC] {TGCT};
            \node (GCTA) [right=of TGCT] {GCTA};
            \node (CTAG) [right=of GCTA] {CTAG};
            \node (TAGC) [right=of CTAG] {TAGC};
            \node (AGCA) [right=of TAGC] {AGCA};

            \node (TGCA) [below=3.6cm of ATGC] {TGCA};
            \node (GCAT) [right=of TGCA] {GCAT};
            \node (CATG) [right=of GCAT] {CATG};
            \node (END) [below=3.6cm of AGCA] {GCA\$};

            % Edges (5-mers)
            \draw[->] (START) -- node[edge label]{\$ATGC} (ATGC);
            \draw[->] (ATGC) -- node[edge label]{ATGCT} (TGCT);
            \draw[->] (TGCT) -- node[edge label]{TGCTA} (GCTA);
            \draw[->] (GCTA) -- node[edge label]{GCTAG} (CTAG);
            \draw[->] (CTAG) -- node[edge label]{CTAGC} (TAGC);
            \draw[->] (TAGC) -- node[edge label]{TAGCA} (AGCA);
            \draw[->] (AGCA) -- node[edge label]{AGCA\$} (END);

            \draw[->] (ATGC) to[] node[edge label,below]{ATGCA} (TGCA);
            \draw[->] (TGCA) -- node[edge label]{TGCAT} (GCAT);
            \draw[->] (GCAT) -- node[edge label]{GCATG} (CATG);
            \draw[->] (CATG) to[] node[edge label,below]{CATGC} (ATGC);

            % Start/end highlights
            \node[draw=none, font=\small, anchor=north] at (START.south) {start};
            \node[draw=none, font=\small, anchor=north] at (END.south) {end};

        \end{tikzpicture}
    }

    \smallskip
    Figure: de~Bruijn graph with 4-mer nodes and 5-mer edge labels.
\end{center}
